% !TeX root = ../../main.tex
\section{Background}
\label{sec:background}

% Review layout: each changed or new paragraph is grouped between "=====" lines.
% For a changed paragraph, the old version is commented out above the new one.
% Paragraphs without a group are unchanged.

The basic DVS model describes each pixel as a detector of changes in log intensity. A positive change reaching the ON threshold produces an ON event, while a negative change reaching the OFF threshold produces an OFF event. After an event, the pixel updates its reference level, and a refractory period limits how soon it can respond again~\cite{lichtsteiner2008dvs,gallego2022event}. This model connects changes in the visual stimulus to the events that a simulator must generate.

In a real camera, however, the effective contrast thresholds are not generally specified directly by the user. Instead, sensor-specific bias settings control circuit operating conditions, including sensitivity, bandwidth, and refractory behavior. Their effects depend on the pixel circuitry and illumination, so a numerical bias setting cannot simply be treated as a contrast threshold with a known physical meaning~\cite{graca2023biasing}. Matching nominal settings is therefore not sufficient to establish that a simulator and a real sensor have the same response.

Further differences arise because pixels do not respond identically or instantaneously. Threshold mismatch changes which brightness variations produce events, finite bandwidth and latency affect their timing, and noise can generate events without the intended stimulus change. These effects also depend on illumination~\cite{hu2021v2e,graca2021noise}. A useful comparison should therefore be sensitive to how events are distributed in space and time, rather than relying only on an aggregate such as the total event count.

% ===== NEW paragraph (outline: existing simulators, moved here from the introduction) =====
Existing simulators reproduce these effects to different degrees, and each focuses on a different part of the process. v2e converts video frames into events with a pixel model that includes intensity-dependent bandwidth, threshold mismatch, and noise, and it interpolates additional frames with a neural network, SuperSloMo, to place events between the original frames~\cite{hu2021v2e}. In our setup, generating the events for a 10~s recording took about 50 minutes, most of it spent on this interpolation. The resulting events still remain grouped around the interpolated frame times, which later work describes as layering along the time axis~\cite{zhang2024v2ce,han2024pecs}. V2CE predicts events from consecutive frames with a learned model and infers continuous timestamps to remove this layering~\cite{zhang2024v2ce}. DVS-Voltmeter models the pixel voltage as a stochastic process that includes the randomness of photon arrival and noise from temperature and parasitic photocurrent~\cite{lin2022dvsvoltmeter}, while PECS starts from a 3D scene and adds the lens and the multispectral rendering of the light reaching the sensor~\cite{han2024pecs}.
% ===== end of new paragraph =====

% ===== Paragraph 4: CHANGED (outline: new opening that links the simulators to the point-set view; the rest is kept) =====
% --- old ---
% Event streams naturally support such a comparison when represented as point sets. Each event provides a spatial location and a timestamp, with polarity as an additional attribute. After alignment and a choice of spatial and temporal scaling, the coordinates can be treated as points in a common space. In this study, we pool polarities and compare three-dimensional space--time points. This representation measures the arrangement of events directly, without requiring features learned for a downstream task, although it does not capture differences in polarity assignment.
% --- new ---
These simulators each address different known limitations, but their improvements are mostly judged by downstream task performance rather than by a direct comparison with real events. A direct comparison needs a representation of the event stream that does not depend on a particular application, and event streams naturally support this when represented as point sets. Each event provides a spatial location and a timestamp, with polarity as an additional attribute. After alignment and a choice of spatial and temporal scaling, the coordinates can be treated as points in a common space. In this study, we pool polarities and compare three-dimensional space--time points. This representation measures the arrangement of events directly, without requiring features learned for a downstream task, although it does not capture differences in polarity assignment.
% ===== end of paragraph 4 =====

% ===== Paragraph 5: CHANGED (one sentence: the evaluation no longer has a sampling reference) =====
% --- old ---
% With this representation, comparing real and simulated recordings becomes a comparison between two empirical distributions of event locations in space and time. The aim is to measure how these distributions differ under the chosen representation. A nonzero value does not by itself establish a simulator error: finite samples from the same distribution can also differ, and repeated real recordings need not be identical. This motivates the sampling and real--real references used in the evaluation. Because the distributions are normalized, event rate must also be considered separately.
% --- new ---
With this representation, comparing real and simulated recordings becomes a comparison between two empirical distributions of event locations in space and time. The aim is to measure how these distributions differ under the chosen representation. A nonzero value does not by itself establish a simulator error: finite samples from the same distribution can also differ, and repeated real recordings need not be identical. This motivates the comparison between repeated real rotations, which serves as a noise floor in the evaluation. Because the distributions are normalized, event rate must also be considered separately.
% ===== end of paragraph 5 =====

MMD measures the difference between two distributions through a kernel, which assigns a similarity to each pair of points~\cite{gretton2012kernel}. It compares the distributions' average representations in the feature space defined by that kernel. With a characteristic kernel, such as a Gaussian kernel, the population MMD is zero only when the distributions are equal. In a finite sample, however, the measured response depends on both sampling and the kernel bandwidth. The bandwidth determines the scale at which differences between event points contribute to the measurement.

Wasserstein distance provides another way to compare distributions by measuring the cost of transporting probability mass from one to the other. The movement cost is defined by distance in the chosen space. SWD uses projections onto one-dimensional axes, computes a Wasserstein distance along each projection, and combines these values~\cite{bonneel2015sliced}. This makes the comparison more tractable for large point sets. Our implementation uses the Python Optimal Transport library~\cite{flamary2021pot}; the projection settings and distance order are specified in the methodology.

% ===== Paragraph 8: CHANGED (outline: Chamfer is used by PECS and kept as an independent verification, not a baseline) =====
% --- old ---
% Chamfer distance provides a baseline based on nearest neighbors, a common approach to comparing point sets~\cite{fan2017pointset}. For each point in one set, it finds the nearest point in the other set, and repeats the comparison in the opposite direction. We use the sum of the two mean Euclidean distances. Unlike Wasserstein transport, this procedure allows many points to share the same nearest neighbor without a constraint on the mass assigned to that neighbor. It therefore provides a different view of the spatial and temporal differences between event sets.
% --- new ---
Chamfer distance compares point sets through nearest neighbors and is commonly used for point clouds~\cite{fan2017pointset}. PECS used it to compare simulated and real events~\cite{han2024pecs}. For each point in one set, it finds the nearest point in the other set, and repeats the comparison in the opposite direction. We use the sum of the two mean Euclidean distances. Unlike Wasserstein transport, this procedure allows many points to share the same nearest neighbor without a constraint on the mass assigned to that neighbor. Because each event is compared only with its closest event in the other set, the value changes with the number of events as well as with their distribution. We therefore include it as an independent verification rather than as one of the main metrics.
% ===== end of paragraph 8 =====

Spike-train distances offer a related approach because events, like neural spikes, occur at discrete times. The van Rossum distance compares temporally filtered spike trains~\cite{vanrossum2001novel}, while SPIKE-distance measures differences in their timing structure~\cite{kreuz2013monitoring}. Applying these methods to event-camera data would also require a rule for handling spatial relationships. Pooling events into a single train discards their locations, whereas comparing corresponding pixels separately does not directly express how a small spatial displacement should affect similarity. We therefore focus on methods that operate jointly on spatial and temporal coordinates.

% ===== NEW paragraph (outline: closing bullet, what the measurement should respond to) =====
These methods provide different ways to compare event distributions, but how their distances relate to physical differences in our recordings still needs to be examined experimentally. We want the measurement to reflect spatial and temporal differences and background noise that change the event distribution, while remaining relatively insensitive to changes in event density that leave the underlying distribution unchanged.
% ===== end of new paragraph =====
